% Generated draft: requires mathematical and source review.
\documentclass[11pt]{article}
\usepackage[margin=25mm]{geometry}
\usepackage{fontspec}
\setmainfont{Latin Modern Roman}
\usepackage{amsmath,amssymb,mathtools,unicode-math}
\setmathfont{Latin Modern Math}
\usepackage{enumitem,needspace,xurl,hyperref}
\setlength{\parindent}{0pt}
\setlength{\parskip}{5pt}
\setlength{\emergencystretch}{3em}
\newcommand{\N}{\mathbb{N}}
\newcommand{\Z}{\mathbb{Z}}
\newcommand{\Q}{\mathbb{Q}}
\newcommand{\R}{\mathbb{R}}
\newcommand{\C}{\mathbb{C}}
\newcommand{\F}{\mathbb{F}}
\newcommand{\A}{\mathbb{A}}
\newcommand{\PP}{\mathbb{P}}
\newcommand{\E}{\mathbb{E}}
\newcommand{\eps}{\varepsilon}
\newcommand{\dd}{\,\mathrm{d}}
\begin{document}
\section*{UnsolvedMath ID: 30000069}
\textbf{Source problem code:} catalogue entry
\subsection{Definitions and mathematical statement}
\textbf{Extracted question.} For increasing events \(A,B,C\) in a percolation model, under what assumptions does\[\mathbb P(A\mid B\cap C)\ge\mathbb P(A\mid C)\] hold? The source specifically asks whether a product-measure assumption or a disconnection-type conditioning event is required. These alternatives are questions, not asserted necessary or sufficient conditions or an exhaustive classification.\par In the usual bond-configuration model \(\Omega=\{0,1\}^{E}\), the coordinate order declares more open edges to be larger. An event \(A\) is increasing when \(\omega\in A\) and \(\omega\le\omega'\) imply \(\omega'\in A\). For a probability measure \(\mathbb P\), conditional probabilities in the display require \(\mathbb P(B\cap C)>0\), which also implies \(\mathbb P(C)>0\). Under these positivity conditions the displayed inequality is equivalent to\[\mathbb P(A\cap B\cap C)\,\mathbb P(C)\ge\mathbb P(A\cap C)\,\mathbb P(B\cap C).\] Thus, for a fixed \(C\), requiring it for every pair of increasing \(A,B\) is positive association of the conditional measure \(\mathbb P(\,\cdot\mid C)\). A product measure means independent edge states, conventionally \(\bigotimes_{e\in E}\operatorname{Bernoulli}(p_e)\); the source does not specify its parameters or whether the model is finite or infinite. Unconditional positive association must not silently be assumed to persist after the unspecified conditioning.\par A nontrivial disconnection event is normally decreasing in this open-edge order, which creates a substantive tension with the wording that all three events are increasing. The source does not identify which conditioning event is intended, the graph, the vertices or sets whose connection is tested, or the class of admissible measures. The displayed formulation is therefore preserved with this ambiguity rather than converted into a particular percolation correlation theorem.\par The supplied triage dated 2026-08-21 labels the item open but says that the extracted inequality and underlying events are missing, preventing identification as a product-measure, BK-type, or later solved variant. The curated statement does display an inequality, so that triage warning and the later \texttt{corrected\_verified} wording are not fully reconciled in the record. No BK-type formulation or present-day status is inferred from them.
\subsection{Short English statement}
Determine when conditioning preserves positive correlation between the specified increasing percolation events. The intended event classes are missing, and the reference to disconnection conditioning conflicts with the usual increasing-event convention. The supplied source-verification and literature notes also require reconciliation.
\subsection{Sources}
\begin{itemize}\item Local UnsolvedMath registry, permanent ID 30000069, code catalogue entry; \emph{Conditional Correlations in Percolation Models}, Oberwolfach Reports Open Problems.\item \emph{Phasen\"uberg\"ange (2005)}, Oberwolfach Reports, DOI 10.4171/owr/2004/30, pp. 1587--1656. \url{https://doi.org/10.4171/owr/2004/30}. Proposed year and year tag: 2004; no proposer named.\item Supplied literature label \emph{original Oberwolfach Report 30/2004}: \url{https://doi.org/10.4171/owr/2004/30}.\item Literature triage checked 2026-08-21, status \texttt{open}; record updated 2026-09-16. The registry separately labels the displayed statement \texttt{corrected\_verified} and reports checking the official report. Neither that verification claim nor the dated literature assessment was independently checked here.\end{itemize}
\end{document}
